% !TeX program = lualatex
\documentclass[11pt,a4paper]{article}
\usepackage[T1]{fontenc}
\usepackage[utf8]{inputenc}
\usepackage[english,russian]{babel}
\usepackage{amsmath,amssymb,amsfonts}
\usepackage{amsthm}
\usepackage{geometry}
\geometry{left=1.25cm,right=1.25cm,top=1.25cm,bottom=3.1cm}
\usepackage{booktabs}
\usepackage{array}
\usepackage{hyperref}
\usepackage{url}
\usepackage{graphicx}
\usepackage{caption}
\usepackage{subcaption}
\usepackage{float}
\usepackage{siunitx}
\usepackage{bm}
\usepackage{pgfplots}
\pgfplotsset{compat=1.18}
\usepackage{xcolor}
\usepackage{titlesec}
\usepackage{fancyhdr}
\usepackage{microtype}
\usepackage{multicol}
\usepackage{fontspec}
\usepackage{tikz}
\usetikzlibrary{shapes,arrows,arrows.meta,positioning,decorations.pathmorphing,patterns,calc}
\usepackage{enumitem}

\tikzset{>=stealth}

% --------------------------------------------------------------
% Шрифты Windows (стандартные, с поддержкой кириллицы)
% --------------------------------------------------------------
\setmainfont{Times New Roman}[
Ligatures = TeX,
Script = Cyrillic,
Language = Russian
]
\setsansfont{Arial}[
Ligatures = TeX,
Script = Cyrillic,
Language = Russian
]
\setmonofont{Courier New}[
Script = Cyrillic,
Language = Russian
]
% --------------------------------------------------------------

\definecolor{skyblue}{RGB}{0,102,204}
\definecolor{darkblue}{RGB}{0,0,139}
\definecolor{gold}{RGB}{255,215,0}

\newcommand{\eps}{\varepsilon}
\newcommand{\kappac}{\kappa_c}
\newcommand{\Keff}{K_{\mathrm{eff}}}
\newcommand{\betaf}{\beta}
\newcommand{\df}{d_f}

\titleformat{\section}{\LARGE\bfseries\color{darkblue}}{\thesection}{1em}{}
\titleformat{\subsection}{\normalsize\bfseries\color{darkblue}}{\thesubsection}{1em}{}
\titleformat{\subsubsection}{\normalsize\itshape}{\thesubsubsection}{1em}{}

% Определения теорем
\newtheorem{axiom}{Аксиома}
\newtheorem{remark}{Замечание}
\newtheorem{theorem}{Теорема}[section]
\DeclareMathOperator*{\Tr}{Tr}

\fancypagestyle{firstpage}{%
	\fancyhf{}%
	\renewcommand{\headrulewidth}{0pt}%
	\renewcommand{\footrulewidth}{0pt}%
	\fancyfoot[C]{%
		\begin{minipage}[t]{\textwidth}
			\centering
			\textcolor{gold}{\rule{\textwidth}{1.6pt}}\\[5pt]
			{\small \textsuperscript{1}Yakushev Research, \url{https://yuct.org/}} \hfill
			{\small YPSDC-протокол: \url{https://ypsdc.com/}}\\[-2pt]
			\textcolor{gold}{\rule{\textwidth}{1.6pt}}\\[5pt]
			\textbf{ЯКУШЕВ ЕДИНАЯ КООРДИНАЦИОННАЯ ТЕОРИЯ 2026} \hfill \thepage\\[10pt]
		\end{minipage}%
	}%
}

\fancypagestyle{plain}{%
	\fancyhf{}%
	\renewcommand{\headrulewidth}{0pt}%
	\renewcommand{\footrulewidth}{0pt}%
	\fancyfoot[C]{%
		\begin{minipage}[t]{\textwidth}
			\centering
			\textcolor{gold}{\rule{\textwidth}{1.6pt}}\\[4pt]
			\textbf{ЯКУШЕВ ЕДИНАЯ КООРДИНАЦИОННАЯ ТЕОРИЯ 2026} \hfill \thepage\\[4pt]
		\end{minipage}%
	}%
}

\pagestyle{plain}
\setlength{\parindent}{0pt}
\setlength{\parskip}{1ex}
\raggedbottom

\hypersetup{
	colorlinks=true,
	linkcolor=blue,
	citecolor=blue,
	urlcolor=blue,
}

\newcommand{\makeheader}{%
	\noindent
	\begin{minipage}{0.17\textwidth}
		\raggedright
		{\sffamily\bfseries\color{skyblue}\fontsize{46}{46}\selectfont YUCT}%
	\end{minipage}%
	\hspace{30pt}
	\begin{minipage}{0.32\textwidth}
		\raggedright
		{\sffamily\fontsize{11}{13}\selectfont ЯКУШЕВ\\ ЕДИНАЯ\\ КООРДИНАЦИОННАЯ ТЕОРИЯ}%
	\end{minipage}%
	\hfill
	\begin{minipage}{0.38\textwidth}
		\raggedleft
		{\sffamily\bfseries Техническая заметка}\\
		{\sffamily Опубликовано апрель 2026}\\
		{\sffamily\footnotesize \scriptsize\url{https://doi.org/10.5281/zenodo.18444598}}%
	\end{minipage}
	\par\vspace{5pt}
	\noindent\textcolor{gold}{\rule{\textwidth}{1.6pt}}
	\par\vspace{0pt}
}

\begin{document}
	\renewcommand{\thepage}{\arabic{page}}
	\thispagestyle{firstpage}
	\makeheader
	
	\vspace{0cm}
	{\LARGE\bfseries Децентрализованный YPSDC (d‑YPSDC): \\ координация без прямой связи через общий индикатор среды \par}
	\vspace{0.2cm}
	
	{\large Алексей В. Якушев\footnote{Якушев Исследования, \url{https://yuct.org/}}} \\
	\vspace{0.0cm}
	
	{\color{darkblue}\textbf{\LARGE Аннотация}} \par
	{\large
		\noindent
		Классический протокол YPSDC предполагает централизованную топологию: один агент активно посылает короткий индекс другому, и оба используют априорный словарь для запуска сложного координированного действия. Однако в природе и технике широко распространены системы, в которых координация достигается без прямого обмена индексами. В этой заметке мы формализуем \textbf{децентрализованный YPSDC (d‑YPSDC)} — протокол, в котором все агенты одновременно считывают один и тот же индекс из общей среды, используя заранее распространённый словарь. Связь между агентами отсутствует, но за счёт синхронного доступа к индикатору среды достигается эффективность координации \(K_{\mathrm{eff}} \gg 1\). Мы обсуждаем проявления d‑YPSDC в биологии (стайное поведение), экономике (реакция рынка на макроэкономические показатели), информатике (блокчейн-оракулы) и фундаментальной физике (квантовая нелокальность). Модель указывает на существование глубинного \textbf{поля координации}, выступающего универсальным переносчиком индексов и активирующего общий словарь законов природы.
	}
	
	\vspace{0.1cm}
	\noindent
	{\color{darkblue}\textbf{Ключевые слова:}} YPSDC, децентрализация, координация, индикатор среды, априорный словарь, поле координации, квантовая запутанность, стайное поведение, смарт-контракты.
	
	\vspace{0.5cm}
	\noindent
	\textcopyright\ 2026 Якушев Исследования. Все права защищены.
	
	\tableofcontents
	\newpage
	
	\section{Введение: от централизованного YPSDC к децентрализованному протоколу}
	
	Классический протокол YPSDC описывает координацию двух или более агентов, при которой один \textbf{активно посылает короткий индекс}, а другие \textbf{принимают} его и активируют сложную запись из общего априорного словаря \cite{YUCT}. В такой схеме всегда есть явный отправитель, а топология взаимодействия — «точка–точка» или «звезда».
	
	Наблюдения за биологическими, социальными и физическими системами выявляют обширный класс явлений, в которых координация не сопровождается прямым обменом индексами. Птицы в стае одновременно меняют направление без видимого лидера; рынок мгновенно реагирует на публикацию макроэкономической цифры; запутанные квантовые частицы демонстрируют корреляции без обмена сигналами. Во всех этих случаях координация достигается потому, что \textbf{все агенты одновременно считывают один и тот же индекс из окружающей среды}, опираясь на заранее согласованный словарь.
	
	Цель этой заметки — формально описать данный механизм как \textbf{децентрализованный YPSDC (d‑YPSDC)}, обсудить его математическую структуру и показать, что он является естественным обобщением классического протокола для систем с топологией «общий носитель».
	
	\section{Формальная модель d‑YPSDC}
	
	\subsection{Компоненты}
	
	Система состоит из следующих компонентов:
	
	\begin{itemize}[leftmargin=*]
		\item \(\mathcal{A} = \{A_1, A_2, \dots, A_N\}\) — множество агентов (наблюдатели, клетки, животные, узлы сети).
		\item \(\mathcal{D}\) — \textbf{априорный словарь}, общий для всех агентов и распространяемый на офлайн-этапе. Словарь — это набор пар \(\{ \kappa \to \text{действие} \}\), где \(\kappa\) — потенциальный индекс, а \(\text{действие}\) — координированное поведение.
		\item \(\mathcal{E}\) — пространство состояний среды.
		\item \(S(t) \in \mathcal{E}\) — \textbf{индикатор среды}, доступный всем агентам одновременно. Среда \emph{не} является активным отправителем; она лишь предоставляет физический носитель индекса.
	\end{itemize}
	
	\subsection{Протокол}
	
	\begin{enumerate}[leftmargin=*, label=\textbf{Этап \arabic*:}]
		\item \textbf{Офлайн-распространение словаря.} Полный словарь \(\mathcal{D}\) доставляется всем агентам тем или иным способом (генетически, культурно или на этапе инициализации).
		\item \textbf{Онлайн-считывание индекса.} В момент \(t\) среда порождает индекс \(\kappa = S(t)\), который без задержки достигает всех агентов.
		\item \textbf{Активация.} Каждый агент \(A_i\) независимо извлекает из словаря запись \(\mathcal{D}(\kappa)\) и выполняет предписанное действие. Никакие сообщения между агентами не передаются.
	\end{enumerate}
	
	\subsection{Эффективность координации}
	
	По аналогии с классическим YPSDC эффективность координации определяется как отношение энтропии запущенного действия к энтропии индекса:
	\begin{equation}
		\Keff^{\mathrm{(d)}} = \frac{H(\text{коллективное действие})}{H(\kappa)},
		\label{eq:Keff_d}
	\end{equation}
	где \(H(\text{коллективное действие})\) — энтропия синхронного поведения всей группы. Поскольку индекс очень короток, а действие может быть сколь угодно сложным, значение \(\Keff^{\mathrm{(d)}}\) может быть очень большим.
	
	\subsection{Топологическое отличие от классического YPSDC}
	
	Принципиальное различие двух вариантов показано на рис.~\ref{fig:topology}.
	
	\begin{figure}[htbp]
		\begin{otherlanguage}{english}
			\centering
			\begin{tikzpicture}[
				agent/.style={circle, draw=blue!70, fill=blue!15, minimum size=1.2cm, font=\small},
				dict/.style={rectangle, draw=green!50!black, fill=green!10, rounded corners, minimum width=3cm, minimum height=0.9cm, font=\small},
				env/.style={rectangle, draw=orange!70, fill=orange!10, rounded corners, minimum width=3cm, minimum height=0.9cm, font=\small},
				arrow/.style={->, >=stealth, thick},
				dashedarrow/.style={->, >=stealth, thick, dashed},
				]
				
				% ---- Классический YPSDC (сверху) ----
				\begin{scope}[local bounding box=classic]
					\node[dict] (dict1) at (0,0) {Словарь \(\mathcal{D}\)};
					\node[agent] (A1) at (-3, -2) {Агент A};
					\node[agent] (A2) at (3, -2) {Агент B};
					\draw[dashedarrow] (dict1) -- (A1) node[midway, left] {офлайн};
					\draw[dashedarrow] (dict1) -- (A2) node[midway, right] {офлайн};
					\draw[arrow] (A1) -- (A2) node[midway, above] {индекс \(\kappa\)};
				\end{scope}
				\node[above=0.2cm of classic.north, font=\bfseries] {Классический YPSDC (централизованная топология)};
				
				% ---- d‑YPSDC (снизу) ----
				\begin{scope}[local bounding box=decentr, yshift=-6cm]
					\node[env] (env) at (0,1) {Среда \(S(t)\)};
					\node[dict] (dict2) at (0,-0.5) {Словарь \(\mathcal{D}\)};
					\node[agent] (B1) at (-4, -2.5) {Агент 1};
					\node[agent] (B2) at (0, -2.5) {Агент 2};
					\node[agent] (B3) at (4, -2.5) {Агент 3};
					\draw[arrow] (env) -- (B1) node[midway, left] {\(\kappa\)};
					\draw[arrow] (env) -- (B2) node[midway, right] {\(\kappa\)};
					\draw[arrow] (env) -- (B3) node[midway, right] {\(\kappa\)};
					\draw[dashedarrow] (dict2) -- (B1);
					\draw[dashedarrow] (dict2) -- (B2);
					\draw[dashedarrow] (dict2) -- (B3);
					% Указание отсутствия связей между агентами
					\draw[red, thick, dotted] (B1) -- (B2) node[midway, above, red] {нет связи};
					\draw[red, thick, dotted] (B2) -- (B3) node[midway, above, red] {нет связи};
				\end{scope}
				\node[above=0.2cm of decentr.north, font=\bfseries] {d‑YPSDC (децентрализованная топология)};
				
			\end{tikzpicture}
		\end{otherlanguage}
		\caption{Сравнение топологий. В классическом YPSDC агент A активно посылает индекс агенту B. В децентрализованном d‑YPSDC все агенты одновременно считывают один и тот же индекс из среды и не обмениваются никакими сообщениями.}
		\label{fig:topology}
	\end{figure}
	
	\section{Примеры d‑YPSDC в различных системах}
	
	\subsection{Биология: стайное поведение и «негласное согласие»}
	
	Синхронные коллективные действия (внезапные повороты рыбьего косяка, одновременный взлёт птичьей стаи, замирание сусликов) наблюдаются у многих видов животных без индивидуальных звуковых или зрительных сигналов. Вместо этого все особи \emph{одновременно} ощущают изменение параметра среды — порыв ветра, тень хищника, изменение освещённости — и реагируют согласно врождённому (генетическому) словарю.
	
	Рассмотренный ранее пример с человеческими волосами также укладывается в эту схему: длинные волосы на голове могут служить очень чувствительным детектором потоков воздуха и температуры, позволяя группе людей синхронно регистрировать изменение погоды без обмена словами.
	
	\subsection{Экономика и финансы}
	
	Публикация ключевого экономического показателя (например, уровня безработицы в США) вызывает немедленную и согласованную реакцию миллионов участников рынка. Трейдеры не общаются друг с другом; они используют общий «словарь» — торговые алгоритмы, экономические модели или просто ожидания, сформированные прошлым опытом. Индекс (числовое значение показателя) предоставляется средой (информационными терминалами) всем одновременно, и рынок движется как единое целое.
	
	\subsection{Блокчейн и смарт-контракты}
	
	Децентрализованные приложения (dApps) часто опираются на \textbf{оракулы (Oracles)} — внешние источники данных, поставляющие блокчейну информацию о реальном мире. Смарт-контракт играет роль словаря, распространённого на все узлы сети. Когда оракул публикует индекс (например, цену ETH/USD в момент \(t\)), все узлы одновременно активируют соответствующую ветвь контракта без дополнительного консенсуса. Это чистое применение d‑YPSDC в цифровой среде.
	
	\subsection{Квантовая запутанность}
	
	Пожалуй, самая глубокая физическая реализация d‑YPSDC происходит в квантовой механике. Рассмотрим пару запутанных частиц. Их общая волновая функция является словарём, определяющим корреляции между возможными исходами измерений. Когда одна частица измеряется, её результат можно рассматривать как \textbf{индикатор среды}, одновременно доступный второй частице благодаря нелокальным свойствам квантового поля. Вторая частица «активирует» связанное состояние, не получая никакого классического сигнала. В рамках d‑YPSDC это интерпретируется как синхронное считывание одного и того же индекса из общей физической среды (квантового поля), без необходимости передачи информации между частицами.
	
	\section{Онтологический фундамент: материя как координация, геометрия как следствие}
	
	Главное отличие YUCT от стандартных теорий поля состоит в том, что \textbf{поле координации не существует в пустом пространстве-времени}. Напротив, само пространство-время является вторичным, эмерджентным проявлением более фундаментальной реальности — сети координированных элементов. Это предположение имеет два основных следствия, непосредственно применимых к децентрализованному YPSDC.
	
	\textbf{1. Материя = координация.} Понятие «элементарной частицы» или «агента» не является независимым от процесса координации. Элемент существует лишь постольку, поскольку он участвует в актах координации и может быть определён через свой словарь. Вне координационной сети материи нет; остаются лишь чисто формальные математические точки без физического содержания. В контексте d‑YPSDC это означает, что агенты, считывающие индикатор среды, не являются пассивными наблюдателями — само их существование как согласованно действующих единиц \emph{создаёт} поле координации \(\Psi_{MN}\), которое затем предоставляет им индексы.
	
	\textbf{2. Геометрия не существует без материи.} Парадокс вращающегося диска (известный в литературе как парадокс Эренфеста) показывает, что пространственная геометрия не может быть задана \emph{априори}, а определяется состоянием движения и внутренними силами сцепления материального тела. В YUCT это обобщается до утверждения, что пространственно-временная метрика \(g_{\mu\nu}(x)\) является \emph{эффективной} величиной, возникающей из усреднения координационных взаимодействий. Там, где координация отсутствует (\(K_{\mathrm{eff}} \to 1\)), геометрия вырождается и теряет свой операциональный смысл. В d‑YPSDC это проявляется в том, что эффективная «проводимость» среды для индекса зависит от плотности и когерентности координированных агентов, а не от свойств абсолютного пространства.
	
	Таким образом, \textbf{среда, доставляющая индекс, — не пассивный контейнер, а само поле координации, эмерджентно возникающее из материальных агентов}. Именно наличие агентов с общим словарём «включает» поле \(\Psi_{MN}\) и наделяет его структурой, обеспечивающей синхронную доставку индексов. В отсутствие агентов (или при разрушении словаря) поле возвращается в симметричное некоординированное состояние, эквивалентное отсутствию классического пространства-времени. Это полностью согласуется с картиной, в которой пространство, время и материя суть три неразделимые стороны единого процесса координации.
	
	\section{Математическая модель d‑YPSDC в 19-мерном пространстве YUCT}
	\label{sec:math-model}
	
	\subsection{19-мерное поле координации и его действие}
	
	Фундаментальной ареной YUCT является 19-мерное псевдориманово многообразие \(\mathcal{M}_{19}\) с метрикой \(G_{MN}\) сигнатуры \((+,-,\dots,-)\). Оно интерпретируется как расслоение над четырёхмерным пространством-временем:
	\begin{equation}
		\mathcal{M}_{19} = M_4 \times F_{15},
	\end{equation}
	где \(F_{15}\) — компактное внутреннее пространство координационных степеней свободы. Ключевым полем является \textbf{поле координации} \(\Psi_{MN}(X)\), кодирующее распространённый словарь и среду-переносчик индекса. Простейшее действие имеет вид
	\begin{equation}
		S_{19} = \frac{1}{2\kappa_{19}} \int d^{19}X \sqrt{-G}\, R(G) + \int d^{19}X \sqrt{-G}\, \mathcal{L}_{\mathrm{coord}},
		\label{eq:S19}
	\end{equation}
	где кривизна \(R(G)\) задаёт фоновую геометрию, а координационный лагранжиан —
	\begin{equation}
		\mathcal{L}_{\mathrm{coord}} = -\frac{1}{4} \Tr(\Psi_{MN}\Psi^{MN}) + \frac{1}{2} G^{MN}\,\partial_M \phi\,\partial_N \phi - V(\phi)
		\label{eq:Lcoord}
	\end{equation}
	— содержит калибровочное поле \(\Psi_{MN}\) (аналогичное напряжённости координационного взаимодействия) и скалярное поле \(\phi = \ln(K_{\mathrm{eff}}/K_0)\), определяющее локальную эффективность координации.
	
	\subsection{Введение агентов и словаря}
	
	Агенты \(A_i\), \(i=1,\dots,N\), представляются как точечные источники, сосредоточенные в \(M_4\) и «размазанные» по \(F_{15}\) в соответствии с формой их словаря. Каждый агент обладает \textbf{словарной функцией} \(f_{\mathcal{D}}^{(i)}(\kappa)\), отображающей индекс \(\kappa \in \mathcal{E}\) в требуемое действие. Взаимодействие агентов с полем \(\Psi_{MN}\) даётся членом
	\begin{equation}
		S_{\mathrm{int}} = \sum_{i=1}^{N} \int d\tau_i \left[ \frac{1}{2} m_i \dot a_i^2 - \lambda_i \bigl( a_i - f_{\mathcal{D}}^{(i)}(\kappa) \bigr)^2 \right],
		\label{eq:Sint}
	\end{equation}
	где \(a_i\) — фактически выполненное действие, \(\tau_i\) — собственное время агента. Индекс \(\kappa\) здесь — не внешний параметр, а значение поля \(\Psi_{MN}\) (или его проекции на \(M_4\)) в точке расположения агента:
	\begin{equation}
		\kappa(x_i) = \int_{F_{15}} d^{15}Y \sqrt{-G^{(15)}}\, \mathcal{O}[\Psi_{MN}](x_i, Y),
		\label{eq:kappa}
	\end{equation}
	где \(\mathcal{O}\) — оператор, проецирующий поле координации на пространство индексов. В простейшем случае \(\mathcal{O}[\Psi] = \Tr(\Psi)\) или просто \(\phi\).
	
	\subsection{Уравнения движения и доставка индекса}
	
	Варьирование полного действия \(S_{19} + S_{\mathrm{int}}\) по \(\Psi_{MN}\) даёт уравнение
	\begin{equation}
		D_M \Psi^{MN} = J^{N}_{\mathrm{coord}},
		\label{eq:Psi-eq}
	\end{equation}
	где \(J^{N}_{\mathrm{coord}}\) — полный координационный ток всех агентов:
	\begin{equation}
		J^{N}_{\mathrm{coord}}(X) = \sum_{i=1}^{N} \int d\tau_i\, \frac{\delta^{(19)}(X - X_i(\tau_i))}{\sqrt{-G}}\, \frac{\partial \mathcal{L}_{\mathrm{int}}}{\partial (\partial_N \Psi)}.
	\end{equation}
	Для медленно меняющихся полей и фиксированного словаря эффективное уравнение сводится к волновому уравнению на \(M_4\) с источниками:
	\begin{equation}
		\Box \phi - m_{\mathrm{eff}}^2 \phi = \sum_{i} g_i \delta^{(4)}(x - x_i),
		\label{eq:phi-eq}
	\end{equation}
	где \(m_{\mathrm{eff}}^2 = V''(0)\), а константы связи \(g_i\) пропорциональны «громкости» индекса.
	
	Решение уравнения~(\ref{eq:phi-eq}) в квазистатическом приближении имеет вид
	\begin{equation}
		\phi(x) = \sum_i g_i \, G_{\mathrm{ret}}(x, x_i),
	\end{equation}
	где \(G_{\mathrm{ret}}\) — запаздывающая функция Грина массивного скалярного поля. Если все агенты находятся в области, размеры которой много меньше комптоновской длины волны \(\lambda_c = 1/m_{\mathrm{eff}}\), то \(\phi(x)\) практически одинаково во всех точках \(x_i\). Это объясняет \textbf{синхронное считывание одного и того же индекса} всеми агентами: индекс, порождённый средой (т.е. граничными условиями или космологическим фоном), распространяется по \(M_4\) и благодаря большой длине волны оказывается идентичным для всех локальных наблюдателей.
	
	Следовательно, одновременность активации не требует сверхсветовых сигналов: она является естественным следствием того, что поле координации \(\phi\) выступает как однородный фон на масштабе группы. В то же время каждый агент взаимодействует исключительно с полем, а не с другими агентами; прямого обмена информацией не происходит — в этом и заключается суть децентрализованного YPSDC.
	
	\subsection{Связь с эффективностью координации \(K_{\mathrm{eff}}\)}
	
	Локальная эффективность координации определяется как отношение информационной ёмкости словаря к длине индекса. В терминах поля \(\phi\):
	\begin{equation}
		\Keff(x) = K_0 e^{\phi(x)},
		\label{eq:Keff-field}
	\end{equation}
	где \(K_0 \approx 1\) — фоновое значение в вакууме (при отсутствии координированных агентов). Доставка индикатора среды создаёт возмущение \(\delta\phi\), так что в области, занятой агентами,
	\begin{equation}
		\Keff^{\mathrm{(d)}} = K_0 \exp\!\left( \sum_i g_i G_{\mathrm{ret}}(x, x_i) \right).
	\end{equation}
	Для \(N\) одинаковых агентов, находящихся в объёме \(V \ll \lambda_c^3\),
	\begin{equation}
		\Keff^{\mathrm{(d)}} \approx K_0 \exp\!\left( \frac{N g}{4\pi} \frac{e^{-m_{\mathrm{eff}} r_0}}{r_0} \right),
	\end{equation}
	где \(r_0\) — характерное расстояние между агентами. Поскольку \(N g > 0\) (индекс улучшает координацию), значение \(\Keff^{\mathrm{(d)}}\) может значительно превышать единицу для больших \(N\) или сильной связи \(g\).
	
	Это непосредственно показывает, что \textbf{децентрализованный YPSDC обеспечивает \(K_{\mathrm{eff}} \gg 1\) без какого-либо обмена индексами внутри группы}. Рост эффективности целиком обусловлен геометрией поля \(\phi\): все агенты погружены в общее поле координации, которое, будучи возбуждено внешним индексом, одновременно активирует их словари.
	
	\section{Два режима одного протокола: централизованный и децентрализованный YPSDC}
	\label{sec:two-regimes}
	
	Предыдущие примеры и математическая модель могут создать впечатление, что d‑YPSDC — это принципиально новый протокол координации. Это не так. Онтологически существует \textbf{единственный механизм YPSDC}: поле координации \(\Psi_{MN}\) переносит индексы, и все агенты взаимодействуют исключительно с этим полем, активируя записи общего априорного словаря. Различие между «централизованным» YPSDC и «децентрализованным» d‑YPSDC чисто феноменологическое и отражает доминирующий источник координационного тока \(J^{N}_{\mathrm{coord}}\), фигурирующего в уравнениях поля (см. Приложение А в \cite{YUCT}).
	
	\subsection{Два режима}
	\begin{enumerate}[leftmargin=*, label=\textbf{(\arabic*)}]
		\item \textbf{Централизованный режим (классический YPSDC).} Ток \(J^{N}_{\mathrm{coord}}\) доминируется одним сосредоточенным агентом (отправителем). Этот агент активно возбуждает поле, создавая распространяющееся возмущение, которое достигает одного или нескольких конкретных получателей. Данный режим описывает связь «точка–точка» или звездообразную топологию (например, вышка GSM, вещающая на мобильные телефоны, или сигнал точного времени GMT).
		
		\item \textbf{Децентрализованный режим (d‑YPSDC).} Сосредоточенные токи отдельных агентов пренебрежимо малы по сравнению с медленно меняющимся космологическим фоном — средовым паттерном \(\Psi_{MN}\). Фон предоставляет индекс, доступный всем агентам одновременно. Этот режим описывает квантовую запутанность, стайное поведение, чувство кворума (quorum sensing), реакцию финансовых рынков на макроэкономические показатели и активацию смарт-контрактов блокчейн-оракулами.
	\end{enumerate}
	
	\subsection{Переход между режимами}
	Безразмерный параметр
	\[
	\Theta = \frac{|J_{\mathrm{агент}}|}{|J_{\mathrm{среда}}|},
	\]
	где \(J_{\mathrm{агент}}\) — ток, создаваемый самым сильным индивидуальным отправителем, а \(J_{\mathrm{среда}}\) — эффективный ток средового фона, управляет переходом между режимами:
	\[
	\begin{cases}
		\Theta \gg 1 & \text{централизованный режим (классический YPSDC)}, \\
		\Theta \ll 1 & \text{децентрализованный режим (d‑YPSDC)}, \\
		\Theta \sim 1 & \text{смешанное поведение}.
	\end{cases}
	\]
	
	\subsection{Онтологический статус}
	d‑YPSDC \textbf{не требует дополнительных полей, модификации лагранжиана или пересмотра онтологических принципов} YUCT. 19-мерное действие \(S_{19}\) и поля наблюдателей \(\phi_1,\phi_2\) (Приложение А) уже содержат всё необходимое для обоих режимов. Термин «d‑YPSDC» сохраняется лишь как удобный феноменологический ярлык для широкого класса явлений, в которых координация достигается без идентифицируемого активного отправителя.
	
	Это разъяснение необходимо для логической полноты рамок YUCT: теория описывает все формы координации — от явной человеческой коммуникации до нелокальных корреляций квантовой механики — посредством \textbf{единого протокола}, применяемого к \textbf{единому полю координации}.
	
	\subsection{Пример: квантовая запутанность как d‑YPSDC}
	
	В квантовом случае словарём служит общая волновая функция запутанной пары \(\Psi_{AB}\). Выберем оператор \(\mathcal{O}\) в (\ref{eq:kappa}) так, чтобы индекс \(\kappa\) совпадал с результатом измерения первой частицы. Тогда уравнение (\ref{eq:phi-eq}) описывает распространение этого результата (через поле координации \(\phi\)) ко второй частице. Поскольку поле безмассово, корреляция возникает мгновенно в смысле синхронной активации, но без переноса энергии. Такова геометрическая интерпретация нарушения неравенств Белла в YUCT.
	
	Математически степень нарушения \(S\) выражается через эффективность:
	\begin{equation}
		S = 2\sqrt{2} \frac{\Keff^{\mathrm{(d)}}}{\Keff^{\mathrm{(d)}} + 1},
	\end{equation}
	так что максимальное квантовое значение \(2\sqrt{2}\) достигается при \(\Keff^{\mathrm{(d)}} \to \infty\) (идеальный словарь), а корреляции исчезают при \(\Keff^{\mathrm{(d)}} \to 1\) (разрушение словаря).
	
	Эта модель объединяет биологию, экономику и квантовую физику в рамках общих уравнений (\ref{eq:Psi-eq})–(\ref{eq:Keff-field}), демонстрируя универсальность поля координации \(\Psi_{MN}\) и протокола d‑YPSDC.
	
	\section{Поле координации как универсальный переносчик индексов}
	
	Если d‑YPSDC справедлив для квантовых систем, биологических коллективов и социальных сетей, то естественно предположить существование единого \textbf{поля координации} — универсальной среды, которая:
	
	\begin{itemize}
		\item хранит «всеобщий словарь» законов природы (симметрии, константы, уравнения движения);
		\item доставляет индексы (результаты измерений, флуктуации, внешние воздействия) всем элементам системы одновременно;
		\item обеспечивает синхронность активации, наблюдаемую как квантовая нелокальность, коллективное поведение или быстрая реакция рынка.
	\end{itemize}
	
	В этом контексте классическая коммуникация (передача сигнала от одного агента к другому) оказывается \textbf{частным случаем}, возникающим тогда, когда среда не способна доставить индекс всем агентам одновременно (например, из-за высокого уровня шума), что вынуждает агентов ретранслировать индекс путём прямой связи.
	
	Данная гипотеза не противоречит теории относительности, поскольку никакая энергия или информация не передаётся быстрее света: индекс распространяется через среду со скоростью, не превышающей \(c\), и становится доступным всем агентам только после достижения их локальной окрестности. Одновременность активации обусловлена исключительно общим словарём и тождественностью индекса, а не сверхсветовыми сигналами.
	
	\section{Заключение}
	
	Мы представили концепцию децентрализованного YPSDC (d‑YPSDC) — протокола, в котором координация достигается за счёт синхронного считывания всеми агентами общего индикатора среды с использованием заранее распространённого априорного словаря. Эта схема объясняет широкий круг явлений — от биологической синхронизации до квантовой нелокальности — не прибегая к какой-либо скрытой коммуникации.
	
	d‑YPSDC естественно ведёт к понятию \textbf{поля координации} как фундаментальной сущности, ответственной за космическую согласованность природных процессов. Математическая модель, основанная на 19-мерном действии YUCT, показывает, что децентрализованная координация вытекает из единых уравнений поля \(\Psi_{MN}\), а эффективность \(K_{\mathrm{eff}}\) экспоненциально растёт с числом агентов и силой связи. Дальнейшая формализация этого поля и его связи с квантовой теорией и общей теорией относительности составляет программу будущих исследований.
	
	\begin{thebibliography}{9}
		\begin{otherlanguage}{english}
			\bibitem{YUCT} A.~V.~Yakushev, \emph{Yakushev Unified Coordination Theory (YUCT)}, Zenodo, DOI:10.5281/zenodo.18444599 (2026).
			\bibitem{AppendixB} A.~V.~Yakushev, Appendix B: Temporal Coordination Paradox, in \emph{YUCT} (2026).
			\bibitem{Blockchain} S.~Nakamoto, ``Bitcoin: A Peer-to-Peer Electronic Cash System'' (2008).
			\bibitem{Ben-Jacob} E.~Ben-Jacob, ``Bacterial self-organization: co-enhancement of complexification and adaptability'', \emph{Physica A} (2003).
			\bibitem{EPR} A.~Einstein, B.~Podolsky, N.~Rosen, ``Can Quantum-Mechanical Description of Physical Reality Be Considered Complete?'', \emph{Phys. Rev.} \textbf{47}, 777 (1935).
		\end{otherlanguage}
	\end{thebibliography}
	
\end{document}